\section{Related work from arXiv, 2025--2026}\label{sec:related}

\subsection{How this list was made, and what can and cannot be trusted}
Search queries on arxiv.org covered: MetaDrive, weather and perception degradation in MetaDrive, robustness benchmarks in MetaDrive, bird's-eye-view representations in MetaDrive, bird's-eye-view perception robustness, and imitation with privileged teachers for driving (searched on 11 November 2026).
For every paper listed below the arXiv abstract page was fetched and the identifier on the page matched the link (HTTP 200); title, authors, first-submission date, problem, method and headline numbers below are taken from that page and \emph{paraphrased} (the abstracts are not reproduced verbatim).
\textbf{The full texts were not read.} Where a statement about the paper's body (for example ``uses MetaDrive'') comes from a search-result summary and not from the abstract page, it is marked \emph{(unverified)}. The one paper outside the 2025--2026 window (the 2021 simulator paper) is marked as such.

\paragraph{The gap this project sits in.} \textbf{None of the abstract pages of the 20 papers below describes a study of weather- or fog/rain-degraded perception with MetaDrive as the simulator.} The closest items are a CARLA fog dataset (FogDrive), an image-perturbation testing framework (PerturbationDrive; simulator not named on the abstract page), and, \emph{(unverified)}, an OpenPilot study in which MetaDrive weather was modelled only as road friction.
Bird's-eye-view driving policies with reinforcement learning exist (ME$^3$-BEV, CARLA) but not in MetaDrive according to the abstract pages. This is a statement about this search, not a proof of novelty: it should be repeated with Google Scholar and the MetaDrive citation list before any novelty claim is written.

\begin{table}[H]\centering\small
\caption{What each paper covers according to its \emph{abstract page}. MD = MetaDrive, W = fog/rain/weather-degraded perception, B = bird's-eye view. ``(?)'' = mentioned only in a search summary of the body text; not verified.}\label{tab:related}
\begin{tabularx}{\textwidth}{@{}lXccc@{}}\toprule
arXiv & Short description (see below for details) & MD & W & B\\\midrule
2109.12674 & MetaDrive simulator (2021, outside the window) & yes & -- & --\\
2512.17586 & Predictive safety representations for safe reinforcement learning & (?) & -- & --\\
2504.18990 & Safety interventions against adversarial patches in an open-source driver assistance system & (?) & (?) friction & --\\
2503.08388 & V-Max: reinforcement-learning framework (Waymax) & -- & -- & --\\
2605.10034 & BehaviorBench: behaviour benchmark for generalisation & (?) & -- & --\\
2503.23708 & Safety and robustness benchmark on safety-critical scenarios & -- & -- & --\\
2605.18059 & Bench2Drive-Robust: robustness under deployment perturbations & -- & -- & --\\
2603.23661 & PerturbationDrive: perturbation testing of driver assistance & -- & yes (mimicked) & --\\
2607.22698 & FogDrive: graded-fog dataset (CARLA) & -- & yes (fog) & --\\
2502.08336 & Salience-invariant visual reinforcement learning & -- & (?) & --\\
2604.04349 & Adversarial robustness of cloud-assisted driving (YOLOv8) & -- & -- & --\\
2508.06074 & ME$^3$-BEV: Mamba + BEV + reinforcement learning (CARLA) & -- & -- & yes\\
2605.30983 & Grace-BEV: graceful degradation of BEV fusion & -- & (?) & yes\\
2603.09529 & RESBev: resilient BEV perception & -- & -- & yes\\
2508.07560 & Survey: progressive BEV perception for safety & -- & yes (adverse weather) & yes\\
2512.20563 & LEAD: learner--expert asymmetry & -- & -- & --\\
2606.19641 & Scaling self-play for end-to-end driving & -- & -- & (prior work only)\\
2505.09546 & Distilling realizable students from unrealizable teachers & -- & -- & --\\
2609.20756 & OPTED: on-policy fine-tuning with a render-free teacher & -- & -- & --\\
2611.12251 & RAIL: representation-aligned imitation learning & -- & -- & --\\\bottomrule
\end{tabularx}
\end{table}

\subsection{MetaDrive and reinforcement-learning frameworks for driving}

\paragraph{MetaDrive: Composing Diverse Driving Scenarios for Generalizable Reinforcement Learning} (Q.~Li, Z.~Peng, L.~Feng et al.; arXiv v1 26 September 2021, v3 16 July 2022; \emph{outside the 2025--2026 window, included because it is the simulator used here}). \url{https://arxiv.org/abs/2109.12674}
\begin{itemize}
\item \textbf{Problem:} generalisation, safe exploration and multi-agent interaction in driving reinforcement learning were studied separately for lack of an integrated environment.
\item \textbf{Method:} a compositional simulator that generates road networks and traffic by procedural generation and imports real-world data; single- and multi-agent tasks with baselines.
\item \textbf{Results:} the abstract states (without numbers) that training on more diverse scenarios improves generalisation. Weather is not mentioned. Our environment configuration (random lane number and width) uses its procedural generation.
\end{itemize}

\paragraph{Learning Safe Autonomous Driving Policies Using Predictive Safety Representations} (M.~Keswani, R.~Bhattacharyya; 19 December 2025). \url{https://arxiv.org/abs/2512.17586}
\begin{itemize}
\item \textbf{Problem:} balancing safety and efficiency in safe reinforcement learning for driving. \textbf{Method:} the SRPL framework (a predictive model of future constraint violations used as a representation) applied to real-world driving data (Waymo Open Motion Dataset, nuPlan), including zero-shot transfer between datasets.
\item \textbf{Results:} statistically significant improvements in success rate (effect sizes $r=0.65$--$0.86$) and in cost reduction ($r=0.70$--$0.83$, $p<0.05$).
\item \textbf{MetaDrive:} not named on the abstract page; a search summary stated that scenarios were built in MetaDrive and that robustness to Gaussian observation noise was tested \emph{(unverified)}. \textbf{Relevance:} safe-RL baselines for a future reinforcement-learning attempt (Section~\ref{sec:rl}).
\end{itemize}

\paragraph{Safety Interventions against Adversarial Patches in an Open-Source Driver Assistance System} (C.~Chen, G.~Xiao, D.~Lee et al.; 26 April 2025). \url{https://arxiv.org/abs/2504.18990}
\begin{itemize}
\item \textbf{Problem:} vulnerability of a driver assistance system to attacks on its perception. \textbf{Method:} simulated safety mechanisms (timely human and automated intervention) evaluated against such attacks. \textbf{Results:} qualitative on the abstract page.
\item \textbf{MetaDrive and weather:} not on the abstract page; a search summary of the body text said the evaluation used MetaDrive and modelled weather only through road friction (for example 50\% friction as heavy rain) \emph{(unverified)}. \textbf{Relevance:} if confirmed, the only item found that links MetaDrive and weather, and it does so through dynamics, which this project deliberately does not do.
\end{itemize}

\paragraph{V-Max: A Reinforcement Learning Framework for Autonomous Driving} (V.~Charraut, W.~Doulazmi, T.~Tournaire et al.; 11 March 2025). \url{https://arxiv.org/abs/2503.08388}
\begin{itemize}
\item \textbf{Problem:} reinforcement learning is rarely adopted for driving because there is no standard, efficient framework. \textbf{Method:} an open framework on the hardware-accelerated Waymax simulator, using ScenarioNet's approach to simulate several datasets. \textbf{Results:} no numbers on the abstract page.
\item \textbf{Relevance:} compares simulators and reports (according to a search summary, \emph{unverified}) that reinforcement-learning agents degrade less than behaviour cloning under noisy initialisation; this agrees qualitatively with our finding that imitation-trained policies are strong in sparse traffic and with the difficulty of speed (a search summary of the body text, \emph{unverified}, says hardware-accelerated simulators are two to three orders of magnitude faster than CPU simulators such as MetaDrive; speed is the bottleneck of our reinforcement-learning attempts).
\end{itemize}

\paragraph{Beyond Self-Play and Scale: A Behavior Benchmark for Generalization in Autonomous Driving} (A.~Distelzweig, F.~Janjo\v{s}, A.~Look et al.; 11 May 2026). \url{https://arxiv.org/abs/2605.10034}
\begin{itemize}
\item \textbf{Problem:} large-scale reinforcement-learning driving policies are not evaluated on standard benchmarks, and existing benchmarks are too simple and use one rule-based traffic model. \textbf{Method:} BehaviorBench: an evaluation interface to nuPlan, an interaction-rich split derived from the Waymo Open Motion Dataset and diverse interactive traffic agents; a hybrid planner combining a PPO policy with a rule-based planner. \textbf{Results:} none on the abstract page.
\item \textbf{Relevance:} supports the evaluation concerns of Direction A. All traffic in our experiments is MetaDrive's own rule-based traffic, a single behaviour model; this is a limitation. MetaDrive is mentioned only in a search summary \emph{(unverified)}.
\end{itemize}

\subsection{Robustness, weather and perturbation benchmarks}

\paragraph{Towards Benchmarking and Assessing the Safety and Robustness of Autonomous Driving on Safety-critical Scenarios} (J.~Li, X.~Liu, S.~Wei et al.; 31 March 2025). \url{https://arxiv.org/abs/2503.23708}
\begin{itemize}
\item \textbf{Problem:} no clear definition of safety-critical driving scenarios and no standard safety evaluation. \textbf{Method:} a definition covering static scenarios (adversarial, distribution shift) and dynamic ones (accidents), and a testing platform for perception and system-level evaluation. \textbf{Results:} none on the abstract page.
\item \textbf{Relevance:} a framework for Directions A and B: the eight early-crash scenes of Section~\ref{sec:failure} are dynamic safety-critical scenarios in this sense.
\end{itemize}

\paragraph{Bench2Drive-Robust: Benchmarking Closed-Loop Autonomous Driving under Deployment Perturbations} (Z.~Zhang, Z.~Jin, Y.~Peng et al.; 18 May 2026). \url{https://arxiv.org/abs/2605.18059}
\begin{itemize}
\item \textbf{Problem:} robustness of closed-loop end-to-end driving to \emph{deployment} imperfections rather than image-level corruption. \textbf{Method:} a closed-loop benchmark with camera-stream failures (frame drop, partial observation), ego-state errors (GPS noise, speed and odometry errors) and compute-induced control delay. \textbf{Results:} the abstract states that such perturbations can substantially degrade driving performance (no numbers).
\item \textbf{Relevance:} Direction E (error in the car's pose against the map is the analogue of GPS offset); our pose-error curve (Figure~\ref{fig:localisation}) is a small instance of the same idea in MetaDrive.
\end{itemize}

\paragraph{PerturbationDrive: A Framework for Perturbation-Based Testing of ADAS} (H.~Leonhard, S.~C.~Lambertenghi, A.~Stocco; 24 March 2026). \url{https://arxiv.org/abs/2603.23661}
\begin{itemize}
\item \textbf{Problem:} deep-network driver assistance is vulnerable to input variations and out-of-distribution data. \textbf{Method:} a tool paper: more than 30 image perturbations that mimic weather, lighting or sensor quality, offline and online closed-loop testing, procedural road generation and search-based testing, ``in different simulators'' (not named on the abstract page). \textbf{Results:} none.
\item \textbf{Relevance:} the closest tool to the teammates' weather module; a candidate for comparing the image-space fog and rain with other perturbation libraries.
\end{itemize}

\paragraph{FogDrive: A Multi-Modal Synthetic Driving Dataset for Perception under Graded Fog} (V.~Panwar, N.~Varshney, P.~W.~Patil; 18 July 2026, v2 9 November 2026). \url{https://arxiv.org/abs/2607.22698}
\begin{itemize}
\item \textbf{Problem:} fog degrades perception in every modality, and real-world capture cannot provide the same scene at known severities. \textbf{Method:} a CARLA dataset of about 660 scenes (about 133\,000 frames), each with a clear reference and three calibrated fog densities, physics-based fog on camera and LiDAR, four cameras, LiDAR, semantic LiDAR and radar.
\item \textbf{Results:} mixing fog densities improves three-dimensional detection at no extra data cost; image-quality metrics (PSNR, SSIM) poorly predict the detection benefit of a dehazer; radar improves camera--LiDAR fusion at every density.
\item \textbf{Relevance:} paired clear/fog frames with graded severity are the same idea as our measurement; the finding that image quality does not predict detection benefit parallels our finding that detection loss, not range error, is what reaches the policy.
\end{itemize}

\paragraph{Salience-Invariant Consistent Policy Learning for Generalization in Visual Reinforcement Learning} (J.~Sun, S.~Tu, Q.~Zhang et al.; 12 February 2025). \url{https://arxiv.org/abs/2502.08336}
\begin{itemize}
\item \textbf{Problem:} visual reinforcement-learning agents overfit to training observations and extract task-irrelevant features in unseen environments. \textbf{Method:} SCPL: a saliency-guided value-consistency module, a dynamics module trained with augmented data and a policy-consistency module with a KL constraint.
\item \textbf{Results:} average improvements of 14\%, 39\% and 69\% on the DeepMind Control video-hard setting, a robotic hard setting and the CARLA benchmark. Weather shift in CARLA (wet, rain) is reported in a search summary \emph{(unverified)}.
\item \textbf{Relevance:} the visual-policy view of the problem; our pixel-policy control (Section~\ref{sec:catalogue}) was meant to show the degradation that such methods address.
\end{itemize}

\paragraph{Adversarial Robustness Analysis of Cloud-Assisted Autonomous Driving Systems} (M.~Al~Islam, A.~S.~El-Wakeel; 6 April 2026). \url{https://arxiv.org/abs/2604.04349}
\begin{itemize}
\item \textbf{Problem:} cross-layer vulnerabilities where perception attacks and vehicle--cloud network impairments together undermine driving. \textbf{Method:} a hardware-in-the-loop testbed with a cloud-hosted YOLOv8 detector attacked by FGSM and PGD, plus network delay (150--250\,ms) and packet loss (0.5--5\%).
\item \textbf{Results:} PGD lowers detection precision and recall from 0.73 and 0.68 (clean) to 0.22 and 0.15 at $\epsilon=0.04$. Not about weather.
\item \textbf{Relevance:} another quantification of how much a YOLOv8 detector can lose; our measurement found a comparable loss (detection 0.72 to 0.15 at rain 0.6) from image-space weather instead of an attack.
\end{itemize}

\subsection{Bird's-eye-view representations and perception}

\paragraph{ME$^3$-BEV: Mamba-Enhanced Deep Reinforcement Learning for End-to-End Autonomous Driving with BEV-Perception} (S.~Lu, R.~Liu, D.~Yang et al.; 8 August 2025). \url{https://arxiv.org/abs/2508.06074}
\begin{itemize}
\item \textbf{Problem:} modular pipelines suffer from error propagation; end-to-end learning is computationally heavy. \textbf{Method:} a Mamba-based spatio-temporal feature network on bird's-eye-view perception feeding an end-to-end deep reinforcement-learning policy. \textbf{Results:} on CARLA it is reported to outperform existing models in collision rate and trajectory accuracy (no numbers on the abstract page).
\item \textbf{Relevance:} bird's-eye view plus reinforcement learning, the combination of our reinforcement-learning attempts; a point of comparison for a stronger reinforcement-learning design (they use a learned spatio-temporal encoder where we use a small network).
\end{itemize}

\paragraph{Can BEV Perception Gracefully Degrade under Sensor Failures?} (H.~Zhang, Y.~Wang, Z.~Zuo; 29 May 2026, v2 5 November 2026). \url{https://arxiv.org/abs/2605.30983}
\begin{itemize}
\item \textbf{Problem:} multi-modal bird's-eye-view fusion models are brittle when a sensor corruption contaminates the shared representation. \textbf{Method:} Grace-BEV, a lightweight robust-fusion framework for Lift-Splat-Shoot detectors with a LiDAR-independent visual expert, a ``TrustGate'' router, a ``FailSafe'' fusion block and three-phase training with modality dropout.
\item \textbf{Results:} on nuScenes-R and nuScenes-C, baselines collapse to 0.0 mAP under complete LiDAR loss while Grace-BEV restores up to 34.7 mAP and improves clean accuracy by up to 1.4 mAP with 1.2\,M extra parameters per detector.
\item \textbf{Relevance:} gives the vocabulary for ``unknown'' versus ``free'' and for degrading gracefully; our grid has an explicit unknown value that the policies do not exploit, a possible extension.
\end{itemize}

\paragraph{RESBev: Making BEV Perception More Robust} (L.~Zhuo, K.~Jin, Z.~Liu et al.; 10 March 2026, v2 30 May 2026). \url{https://arxiv.org/abs/2603.09529}
\begin{itemize}
\item \textbf{Problem:} bird's-eye-view perception suffers from sensor degradation and adversarial attacks. \textbf{Method:} a plug-and-play latent world model that predicts clean bird's-eye-view features to reconstruct corrupted observations, working at the semantic feature level of Lift-Splat-Shoot pipelines, with few-shot fine-tuning. \textbf{Results:} the abstract reports significantly improved robustness on nuScenes without numbers.
\item \textbf{Relevance:} a perception-side counterpart to the policy-side robustness question of this project.
\end{itemize}

\paragraph{Progressive Bird's Eye View Perception for Safety-Critical Autonomous Driving: A Comprehensive Survey} (Y.~Gong, N.~Wang, J.~Lu et al.; 11 August 2025). \url{https://arxiv.org/abs/2508.07560}
\begin{itemize}
\item \textbf{Problem:} whether bird's-eye-view perception is safe and reliable under occlusion, adverse weather and dynamic traffic. \textbf{Method:} a survey of single-modality vehicle-side, multimodal vehicle-side and multi-agent collaborative bird's-eye-view perception, with datasets. \textbf{Results:} none (survey).
\item \textbf{Relevance:} background for the representation choices of Section~\ref{sec:arch}.
\end{itemize}

\subsection{Imitation, privileged teachers and on-policy distillation}

\paragraph{LEAD: Minimizing Learner--Expert Asymmetry in End-to-End Driving} (L.~Nguyen, M.~Fauth, B.~Jaeger et al.; 23 December 2025, v2 13 April 2026; accepted at CVPR 2026). \url{https://arxiv.org/abs/2512.20563}
\begin{itemize}
\item \textbf{Problem:} privileged experts see more and are more certain than sensor-based students, and navigational intent is under-specified, which limits imitation. \textbf{Method:} an empirical study of the asymmetries and practical interventions; the student is the TransFuser version 6 policy (CARLA). \textbf{Results:} 95 driving score on Bench2Drive and more than double the previous performance on Longest6~v2 and Town13.
\item \textbf{Relevance:} \emph{the closest published statement of Direction C}: our teacher sees 45\,m of exact state while the student's grid ends at 17.5\,m.
\end{itemize}

\paragraph{Scaling Self-Play for End-to-End Driving} (L.~Rowe, R.~Girgis, R.~de~Schaetzen et al.; 17 June 2026, v2 19 June 2026). \url{https://arxiv.org/abs/2606.19641}
\begin{itemize}
\item \textbf{Problem:} models trained on offline demonstrations suffer compounding errors and have no closed-loop feedback. \textbf{Method:} large-scale self-play from pixels in ``Gigapixel'', a batched simulator with perspective rendering of a simplified bounding-box world; a reinforcement-learning teacher on vectorised observations is distilled with self-play dataset aggregation, followed by perception adaptation to real sensor data.
\item \textbf{Results:} a throughput of 50\,000 agent steps per second and competitive performance on HUGSIM and NAVSIM version~2.
\item \textbf{Relevance:} validates the teacher--student dataset-aggregation design and shows how far simulator speed can be pushed compared with the CPU simulator used here.
\end{itemize}

\paragraph{Distilling Realizable Students from Unrealizable Teachers} (Y.~Kim, N.~Chin, A.~Vasudev et al.; 14 May 2025). \url{https://arxiv.org/abs/2505.09546}
\begin{itemize}
\item \textbf{Problem:} information asymmetry between a privileged teacher and a student causes distribution shift and degraded policies. \textbf{Method:} an imitation-learning scheme that queries the teacher adaptively and a reinforcement-learning scheme that chooses where to start training from recovery states. \textbf{Results:} significant improvements over standard teacher--student baselines on simulated and real robotic tasks (no numbers on the abstract page).
\item \textbf{Relevance:} describes the failure mode of Direction C and proposes remedies that could be tried here.
\end{itemize}

\paragraph{OPTED: On-Policy Fine-Tuning for End-to-End Driving using a Render-Free Teacher} (D.~Da~Col, M.~Igl, P.~Karkus et al.; 17 September 2026). \url{https://arxiv.org/abs/2609.20756}
\begin{itemize}
\item \textbf{Problem:} pre-trained end-to-end policies drift in closed loop and sensor-based reinforcement-learning post-training is expensive. \textbf{Method:} a teacher trained by reinforcement learning on vectorised inputs (HD map and bounding boxes) supervises a pre-trained camera-based student during closed-loop fine-tuning in AlpaSim, on three-dimensional Gaussian-splatting reconstructions of real logs (TransFuser and VaVAM).
\item \textbf{Results:} driving scores increase by factors of 1.6 and 9.5 for the two models, with about three orders of magnitude fewer simulator interactions than direct reinforcement-learning post-training.
\item \textbf{Relevance:} the strongest argument for distilling from a privileged teacher instead of fine-tuning the student with reinforcement learning, which is what our four attempts found.
\end{itemize}

\paragraph{RAIL: Representation-Aligned Imitation Learning for Student-Compatible Teacher Policies} (M.~Mammadov, P.~Zuidberg~Dos~Martires, J.~A.~Stork; 9 November 2026; accepted at NeurIPS 2026 according to the page's comments field). \url{https://arxiv.org/abs/2611.12251}
\begin{itemize}
\item \textbf{Problem:} when a privileged teacher cannot be imitated exactly from raw inputs an irreducible imitation gap remains. \textbf{Method:} contrastive learning of a latent representation shared by teacher and student observations; the teacher is trained in that space so that it stays imitable. \textbf{Results:} outperforms strong baselines and substantially reduces the gap (no numbers on the abstract page).
\item \textbf{Relevance:} a method-level answer to the information gap of Section~\ref{sec:failure}.
\end{itemize}
